% !TeX root = ../教材题库_全部条目_转写工作稿.tex
% Source: XJ-B1-C04-EX-P131-Y0177
% 第4章 / 4.4 函数与方程 / 教材 p.131 / PDF 第138页
% OCR未校对：公式与文字均可能有误；未完成数学TeX转写；禁止直接发给学生
% 原题及相邻原解截图（核对用）：../题目原图/XJ-B1-C04-EX-P131-Y0177-01.webp
% 本文件只在转写完成后作为题干使用；OCR 条目尚不是可靠题干。
\begin{bankitem}{4.4 例题 例题2}{第4章 / 4.4 函数与方程 / 教材 p.131 / PDF 第138页}{XJ-B1-C04-EX-P131-Y0177}{OCR未校对：公式与文字均可能有误；未完成数学TeX转写；禁止直接发给学生}
\begin{ocrdraft}
[教材p.131 / PDF第138页]
\par
00讨论方程2-x=lo\ocrglyph{00E6}x的解的个数与分布情况．
\par
解
\par
（方法一）在平面直角坐标系内，分别画出函数一2一'与的大致
\par
图象，如图4，2一2所示．
\par
由于函数y一2一'，一log2工分别是减函数
\par
Y'
\par
和增函数，由图可看出，函数3，一2一'与3，=
\par
2
\par
，logzx
\par
logz工的图象只在（0，十。。）上有一个交点，所
\par
以方程2一'一logzx有且只有一个解．
\par
2
\par
3
\par
1
\par
由于1时，2-1=
\par
一，log21一0，于是
\par
2
\par
2-l>log21.
\par
图4．4
\par
一2
\par
1
\par
又工一2时，2一2=
\par
一，10g22一1，于是2一2<Iog22．
\par
4
\par
综上可知，方程只有一个解，且在区间（1，2）上．
\par
（方法二）考察函数不到一logz工一2一'，e（0，十亠），其零点就是方程2一
\par
logz工的解．
\par
1
\par
3
\par
由于
\par
不1）一10g21一2一
\par
一<0，不2）==Iogz2\ocrglyph{00AF}2-
\par
2
\par
4
\par
可见不到一logz工一2一在（1，2）内有零点．
\par
另一方面，由于函数3-lo\ocrglyph{00E6}x单调递增且函数3，一一2一'也单调递增（因为函数
\par
3，一2一'单调递减），因而函数不到一logz工一2一'单调递增，所以丆0）一logz工一2一'在
\par
（1，2）内只有一个零点．因此，方程只有一个解，且在区间（1，2）上．
\end{ocrdraft}
\end{bankitem}
